\setcounter{tocdepth}{1}
\tableofcontents
\subsection*{Notes}
Each topic contains \emph{core problems}, covering the minimum required for that topic;
\emph{additional problems}, providing variety and extending the core material; and
\emph{challenges}, related to the topic but going beyond the requirements of standard
introductory courses in probability and statistics.
The following symbols are used:
\begin{itemize}
  \item[\faBook] a problem introducing fundamental objects of the course;
  \item[\faStar] a challenging problem;
  \item[\faSmile] an interesting example;
  \item[\faPencil*] a problem from an FMI test;
  \item[\faRobot] an independent research task. All resources are permitted, including
  literature, software and AI. You are expected to check the results, understand them,
  and be able to explain them.
\end{itemize}
{\small
\textbf{Note on the use of AI.}
The latest publicly available versions of ChatGPT and Claude (at the time of writing,
sol 6.1 and opus 5.5, respectively) were used for formatting, editing and clarifying
individual formulations. The compiler is responsible for the final version and its
mathematical correctness.
}
\section{Combinatorics}
\subsection*{Core problems}
\begin{problembook}
For $n,k\ge 1$, recall:
\begin{enumeratea}
  \item how many solutions does $x_1+\cdots+x_n=k$ have when $x_1,\dots,x_n$
  are positive integers? What if they are nonnegative integers?
  \item state the inclusion--exclusion principle.
\end{enumeratea}
\end{problembook}
\begin{problembook}
Let $M$ be a set with $n$ elements. In how many ways can we choose:
\begin{enumeratea}
  \item $k$ distinct elements of $M$;
  \item an ordered $k$-tuple of distinct elements of $M$;
  \item an ordered $k$-tuple of elements of $M$, with repetitions allowed?
\end{enumeratea}
\end{problembook}
\begin{problem}
In how many ways can all $k$ distinguishable particles be distributed among $n$
distinct cells if:
\begin{enumeratea}
  \item each cell contains at most one particle;
  \item cells may contain any number of particles;
  \item no cell is empty?
\end{enumeratea}
Answer the same questions for indistinguishable particles. Also state when the answer is zero.
\end{problem}
\begin{problem}
Five distinguishable balls are distributed among three distinct boxes $A,B,C$.
Find the number of distributions in which:
\begin{enumeratea}
  \item $A$ is empty;
  \item only $A$ is empty;
  \item exactly one box is empty;
  \item at least one box is empty;
  \item no box is empty.
\end{enumeratea}
\end{problem}
\begin{problemstar}
Consider the following questions.
\begin{enumeratea}
  \item In how many ways can we go from $(0,0)$ to $(n,n)$ on the square lattice
  using unit steps to the right and upwards? What if the path never goes above the diagonal?
  \item How many correctly balanced parenthesis expressions have $n$ opening and $n$
  closing parentheses? Construct an explicit bijection between them and the paths
  in the second part of the preceding question.
  \item Find simple asymptotic equivalents for the total number of paths and for the
  number that never go above the diagonal, as $n\to\infty$.
\end{enumeratea}
\end{problemstar}
\subsection*{Additional problems}
\begin{problem}
How many four-digit numbers can be formed from the digits $1,2,3,4,5$ if:
\begin{enumeratea}
  \item repetition is not allowed;
  \item repetition is allowed;
  \item repetition is not allowed and the number is odd?
\end{enumeratea}
\end{problem}
\begin{problem}
In how many ways can a four-person delegation be selected from 12 candidates,
including distinct candidates $A,B,C,D$, if:
\begin{enumeratea}
  \item there are no restrictions;
  \item $A$ and $B$ must not both participate;
  \item either both $C$ and $D$ participate, or neither does?
\end{enumeratea}
\end{problem}
\begin{problem}
Prove combinatorially, and then algebraically, that
\[
\sum_{k=0}^n\binom nk=2^n
\qquad\text{and}\qquad
\binom{2n}{n}=\sum_{k=0}^n\binom nk^2.
\]
\end{problem}
\begin{problem}
Let $p(n)$ be the number of partitions of $n$ into distinct positive integers,
and $q(n)$ the number of partitions into odd positive integers, where repetitions
are allowed and the order of the summands does not matter. For example, for $n=6$:
\[
p(6)=4:\ 6,\ 5+1,\ 4+2,\ 3+2+1; \qquad
q(6)=4:\ 5+1,\ 3+3,\ 3+1+1+1,\ 1+1+1+1+1+1.
\]
\begin{center}
\begin{tabular}{c|cccccccc}
  $n$    & 1 & 2 & 3 & 4 & 5 & 6 & 7 & 8 \\ \hline
  $p(n)$ & 1 & 1 & 2 & 2 & 3 & 4 & 5 & 6 \\
  $q(n)$ & 1 & 1 & 2 & 2 & 3 & 4 & 5 & 6
\end{tabular}
\end{center}
By comparing generating functions, prove that $p(n)=q(n)$ for every $n$.
\end{problem}
\subsection*{Challenges and independent research}
\begin{problemstar}[Derangements]
Let $D_n$ be the number of permutations $\sigma\in S_n$ such that $\sigma(i)\ne i$ for every $i$.
\begin{enumeratea}
  \item Prove that
  \[ D_n=n!\sum_{j=0}^{n}\frac{(-1)^j}{j!}. \]
  \item Prove combinatorially that, for $n\ge 2$,
  \[ D_n=(n-1)(D_{n-1}+D_{n-2}), \]
  where $D_0=1$ and $D_1=0$.
  \item Prove that, for $n\ge 1$, $D_n$ is the nearest integer to $n!/e$,
  without using a numerical approximation to $e$.
\end{enumeratea}
\end{problemstar}
\begin{problemstar}[Necklaces and Burnside's lemma]
Let $X=\{0,1\}^n$ be the set of binary words of length $n$, and let
$G=\{r^0,r^1,\dots,r^{n-1}\}\cong\mathbb Z_n$ be the group of cyclic rotations,
$r(x_1x_2\dots x_n)=x_nx_1\dots x_{n-1}$. A \emph{necklace} is an orbit of $G$
in $X$: two words are identified if one is obtained from the other by a cyclic
rotation. For example, $0011,\,0110,\,1100,\,1001$ form one necklace, and
$0101,\,1010$ form another. Rotations preserve the number of ones, so $G$ also acts on
\[ X_k=\{x\in X:\ x\text{ has exactly }k\text{ ones}\}. \]
Write $N(n,k)$ for the number of orbits in $X_k$, that is, the number of necklaces
with $n$ beads and exactly $k$ ones.
For $g\in G$, let
\[ \operatorname{Fix}(g)=\{x\in X_k:\ g(x)=x\}, \]
and let $X_k/G$ be the set of orbits. Use Burnside's (Cauchy--Frobenius) lemma
\[ \#(X_k/G)=\frac1{|G|}\sum_{g\in G}|\operatorname{Fix}(g)|. \]
Explain why a rotation $g$ whose cycles on the positions $\{1,\dots,n\}$ have
length $d$ can fix a word with $k$ ones only if $d\mid k$. Prove that in this case
\[ |\operatorname{Fix}(g)|=\binom{n/d}{k/d}, \]
whereas if $d\nmid k$, then $\operatorname{Fix}(g)=\emptyset$. Also prove that
for every $d\mid n$ there are exactly $\varphi(d)$ rotations whose cycles have length $d$.
Deduce that, for $1\le k\le n-1$,
\[ N(n,k)=\frac1n\sum_{d\mid\gcd(n,k)}\varphi(d)\binom{n/d}{k/d}, \]
where $\varphi$ is Euler's totient function. As a check, find $N(12,4)$.
\end{problemstar}
\begin{researchtask}[Random mappings and cyclic vertices]
Let $f:[n]\to[n]$ be chosen uniformly from all $n^n$ mappings. Let $G_f$ be the
directed graph with edges $i\to f(i)$, and $K_n$ the number of vertices lying on
a cycle in $G_f$. You may use AI to propose solutions, acting as their \emph{reviewer}:
verification, rather than generation, is assessed.
\begin{enumeratea}
  \item Use software to find the distribution $\Pbb(K_n=k)$ for $n\le 7$.
  The code may be AI-generated, but run it and check that
  $\sum_k\#\{K_n=k\}=n^n$; also check the case $n=3$ by hand.
  \item Find and prove a formula for $\#\{f:K_n=k\}$. You may ask AI for a proposal.
  (i) Compare it with the table from part 1. (ii) Write the proof in your own words,
  separating the choice of cyclic vertices, the permutation on them, and the directed
  forest rooted at them. (iii) Give your own proof of the formula $k\,n^{n-k-1}$
  for the number of forests with $k$ prescribed roots. If you use AI, check whether
  it has actually proved the formula, and give a proof by a different method.
  (iv) Find and numerically check the number of mappings with $K_n=k$ and exactly $c$ cycles.
  \item Derive $\Ebb K_n$ in two ways: from the distribution and by linearity,
  computing $\Pbb(\text{vertex }1\text{ is cyclic})$. Prove the identity between
  the two expressions. Use the exact formula to calculate $\Ebb K_n$ for
  $n=10,10^2,10^3,10^4,10^5$ and investigate numerically whether
  \[ \frac{\Ebb K_n}{\sqrt n}\longrightarrow\sqrt{\frac{\pi}{2}}. \]
  Explain the origin of the constant $\sqrt{\pi/2}$. Optionally investigate the finer asymptotic claim
  \[ \Ebb K_n-\sqrt{\frac{\pi n}{2}}\longrightarrow-\frac13. \]
  Let $T_n$ be the number of distinct vertices in the trajectory $1,f(1),f^2(1),\dots$
  before its first repetition. Do $K_n$ and $T_n$ have the same distribution,
  or only the same expectation?
\end{enumeratea}
\end{researchtask}
\noindent\textbf{Lean.} Let $A,B\subseteq\Omega$, with indicators taking values in $\mathbb N$.
Define the indicator by
\[
\ind_A(x)=\begin{cases}1,&x\in A,\\0,&x\notin A,\end{cases}
\]
and prove in Lean 4 with Mathlib that, for every $x\in\Omega$,
\[ \ind_{A\cup B}(x)+\ind_{A\cap B}(x)=\ind_A(x)+\ind_B(x). \]
